%------------------------------------------------------------------------------%
% \chapter{Plano Cartesiano e Equações}
%------------------------------------------------------------------------------%

%------------------------------------------------------------------------------%
\begin{exercicios}{Plano Cartesiano}
    
    \begin{exercicio}
        Posicione os pontos $(3,2)$, $(4,1)$, $(-1,-1)$, $(\sfrac{1}{2},\sfrac{3}{4})$ e $(\pi,\e)$ nos plano cartesiano.
        \begin{answers} \(\qquad\) \addgraphic{.6}{answer-plano-cartesiano-01}    
        \end{answers}
    \end{exercicio}
    
    \begin{exercicio}
        Escreva os pares ordenados correspondentes aos pontos da figura abaixo.
        \addgraphic{.9}{exercicio-plano-cartesiano}
        \begin{answers}
            \begin{mathlist}[3]
                \item A=(1,8)
                \item B=(-8,8)
                \item C=(-5,6)
                \item D=(-8,3)
                \item E=(-6,-2)
                \item F=(2,-3)
                \item G=(6,-5)
                \item H=(7,4)
                \item I=(5,-2)
                \item J=(4,2)
                \item K=(7,7)
                \item L=(-3,-4)
                \item M=(-7,-6)
                \item N=(7,-2)
                \item P=(7,0)
                \item Q=(-1,0)
                \item R=(0,3)
                \item S=(0,-8)
            \end{mathlist}

        \end{answers}
    \end{exercicio}

    \begin{exercicio}
        Represente, no plano Cartesiano, os pontos 
        $(0,4)$, $(1,0)$, $(1,3)$, $(5,1)$, $(-2,0)$, $(0,-3)$, $(3,-4)$, $(4,-2)$, $(-6,2)$, $(-3,-5)$ e $(-4,-1)$.
        \begin{answers}
            \addgraphic{.6}{answer-plano-cartesiano-03}
        \end{answers}
    \end{exercicio}
    
    \begin{exercicio}
        Se o ponto $P(x,y)$ está no segundo quadrante, quais são os sinais de $x$ e $y$?
    \end{exercicio}
    
    \begin{exercicio}
        Exiba no plano Cartesiano as regiões definidas pelos conjuntos abaixo.
	    \begin{mathlist}[2]
	    	\item \{(x, y) \st x > - 1/2\}
	    	\item \{(x, y) \st -3 \leq y \leq 0\}
	    	\item \{(x, y) \st y \geq 1,5\}
	    	\item \{(x, y) \st x = 3\}
	    	\item \{(x, y) \st -2 \leq x \leq 5\}
	    	\item \{(x, y) \st x \leq -1\}
	    	\item \{(x, y) \st y < 3/2\}
	    	\item \{(x, y) \st y = -2\}
	    \end{mathlist}      
    \end{exercicio}
    
    \begin{exercicio}
        Exiba no plano Cartesiano as regiões definidas abaixo.
	    \begin{mathlist}
	    	\item x \geq 1 \text{ e } y \geq 1
	    	\item x \geq - 1 \text{ e } y \leq 2
	    	\item -3 \leq x \leq 2 \text{ e } 0 \leq y \leq 3
	    	\item -4 \leq x \leq 1 \text{ e } y = 1
	    \end{mathlist}        
    \end{exercicio}
    
    \begin{exercicio}
        Expresse o conjunto de pontos do primeiro quadrante usando desigualdades.
    \end{exercicio}

    \begin{exercicio}
        Verifique, por substituição, se os pontos abaixo pertencem ao gráfico da equação correspondente.
	    \begin{alphalist}
	    	\item \( y = x^3 - 4x^2 + 5x - 12 \)           \tabto{46mm} $(4,10)$                        \tabto{60mm} $(-3,-9)$
	    	\item \( y = \sqrt{x^2 + 1} \)                 \tabto{46mm} $(7,5\sqrt{2})$                 \tabto{60mm} $(0,-1)$
	    	\item \( y = |3x - 8| \)                       \tabto{46mm} $\left(-\sfrac{1}{3},7\right)$  \tabto{60mm} $\left(\sfrac{1}{3}, 7\right)$
	    	\item \( y = \dfrac{1}{|x-2|}-\frac{1}{x-2} \) \tabto{46mm} $\left(-\sfrac{5}{2},0\right)$  \tabto{60mm} $\left(-6,-\!\sfrac{1}{8}\right)$
	    	\item \( y = \dfrac{x^2-2x+1}{-x^3+2x^2+4} \)  \tabto{46mm} $\left(1,\sfrac{2}{3}\right)$   \tabto{60mm} $(2, 4)$
	    	\item \( y^2 + xy^2 - 9 = x^2 + 3xy - 13 \)    \tabto{46mm} $(-2, 6)$                       \tabto{60mm} $(5,-1)$
	    \end{alphalist} 
    \end{exercicio}
    
    \begin{exercicio}
        Usando uma tabela de pares $(x,y)$, trace o gráfico das equações abaixo, no intervalo especificado.
        \begin{mathlist}
            \item y = -2x + 3            \) \tabto{30mm} $x \in [-2, 3]$
            \item 3y - 2x + 3 = 0        \) \tabto{30mm} $x \in [-2, 4]$
            \item 2y + x = 4             \) \tabto{30mm} $x \in [-2, 6]$
            \item y = x^2 - 1            \) \tabto{30mm} $x \in [-2, 2]$
            \item y = 2 - \frac{1}{2}x^2 \) \tabto{30mm} $x \in [-3, 3]$
            \item y = -2x^2 + 4x         \) \tabto{30mm} $x \in [-1, 3]$
            \item y = x^3 - x            \) \tabto{30mm} $x \in [-2, 2]$
            \item y = \sqrt{x}           \) \tabto{30mm} $x \in [ 0, 4]$
            \item y = \sqrt{x + 1}       \) \tabto{30mm} $x \in [-1, 3]$
            \item y = 1/x                \) \tabto{30mm} $x \in [-4, 4]$
            \item y = \abs{x}            \) \tabto{30mm} $x \in [-2, 2]$
            \item y = \abs{x - 2}        \) \tabto{30mm} $x \in [-1, 5]$
        \end{mathlist}   
        \begin{answers}
            \begin{alphalist}[2]
                \item \addgraphic{.6}{answer-plano-cartesiano-09-01}
                \item \addgraphic{.6}{answer-plano-cartesiano-09-02}
                \item \addgraphic{.6}{answer-plano-cartesiano-09-03}
                \item \addgraphic{.6}{answer-plano-cartesiano-09-04}
                \item \addgraphic{.6}{answer-plano-cartesiano-09-05}
                \item \addgraphic{.6}{answer-plano-cartesiano-09-06}
                \item \addgraphic{.6}{answer-plano-cartesiano-09-07}
                \item \addgraphic{.6}{answer-plano-cartesiano-09-08}
                \item \addgraphic{.6}{answer-plano-cartesiano-09-09}
                \item \addgraphic{.6}{answer-plano-cartesiano-09-10}
                \item \addgraphic{.6}{answer-plano-cartesiano-09-11}
                \item \addgraphic{.6}{answer-plano-cartesiano-09-12}
            \end{alphalist}
        \end{answers}
    \end{exercicio}
    
    \begin{exercicio}
        Determine os interceptos e trace o gráfico das equações abaixo.
        \begin{mathlist}[2]
            \item y = x - 1
            \item y = x^2 + 2x - 3
            \item y = 3 - x/2
            \item y = -x^2 + 8x - 12
            \item x = 3 + y
            \item x = y^2 - 1
            \item x = (y+1)(y-2)
            \item x + y = 2
        \end{mathlist} 
        \begin{answers}
            \begin{alphalist}[2]
                \item \addgraphic{.6}{answer-plano-cartesiano-10-01}
                \item \addgraphic{.6}{answer-plano-cartesiano-10-02}
                \item \addgraphic{.6}{answer-plano-cartesiano-10-03}
                \item \addgraphic{.6}{answer-plano-cartesiano-10-04}
                \item \addgraphic{.6}{answer-plano-cartesiano-10-05}
                \item \addgraphic{.6}{answer-plano-cartesiano-10-06}
                \item \addgraphic{.6}{answer-plano-cartesiano-10-07}
                \item \addgraphic{.6}{answer-plano-cartesiano-10-08}
            \end{alphalist}
        \end{answers}
    \end{exercicio}     
    
    \begin{exercicio}
        Represente os pontos em um sistema de eixos coordenados.
        \begin{mathlist}[2]
            \item (-2, 5)
            \item (3, -1)
            \item (3, -4)
            \item (8, -7/2)
            \item (4,5; -4,5)
        \end{mathlist}
        \begin{answers} \(\qquad\) \addgraphic{.6}{answer-plano-cartesiano-11}
        \end{answers}
    \end{exercicio}
   
\end{exercicios}

%------------------------------------------------------------------------------%
